\documentclass[11pt]{article}

\usepackage[T1]{fontenc}
\usepackage[utf8]{inputenc}
\usepackage{lmodern}
\usepackage[a4paper,margin=1in]{geometry}
\usepackage{microtype}
\usepackage{amsmath,amssymb,amsthm,mathtools}
\usepackage{mathrsfs}
\usepackage{bm}
\usepackage{enumitem}
\usepackage{booktabs}
\usepackage{xcolor}
\usepackage{longtable}
\usepackage{array}
\usepackage{hyperref}
\usepackage[nameinlink,noabbrev]{cleveref}

\hypersetup{
  colorlinks=true,
  linkcolor=blue!55!black,
  citecolor=blue!55!black,
  urlcolor=blue!55!black,
  pdftitle={Audited SOMK--GFT Reduction for Binary Goldbach},
  pdfauthor={Draft manuscript}
}

\newtheorem{theorem}{Theorem}[section]
\newtheorem{lemma}[theorem]{Lemma}
\newtheorem{proposition}[theorem]{Proposition}
\newtheorem{corollary}[theorem]{Corollary}
\newtheorem{hypothesis}[theorem]{Hypothesis}
\newtheorem{conjecture}[theorem]{Conjecture}
\theoremstyle{definition}
\newtheorem{definition}[theorem]{Definition}
\theoremstyle{remark}
\newtheorem{remark}[theorem]{Remark}
\newtheorem{warning}[theorem]{Warning}

\newcommand{\T}{\mathbb T}
\newcommand{\Pp}{\mathbb P}
\newcommand{\1}{\mathbf 1}
\newcommand{\e}{\mathrm e}
\newcommand{\eps}{\varepsilon}
\newcommand{\Mcal}{\mathfrak M}
\newcommand{\mcal}{\mathfrak m}
\newcommand{\SG}{\mathfrak S}
\newcommand{\Jw}{\mathfrak J_W}
\newcommand{\PMKF}{\mathrm{PMKF}\text{-}\mathrm{GFT}}
\newcommand{\SOMK}{\mathrm{SOMK}\text{-}\mathrm{GFT}}
\newcommand{\PBM}{\mathrm{PB}\text{-}\mathrm{MRDET}}
\newcommand{\Kl}{\mathcal K}
\newcommand{\Eend}{\mathcal E_{\mathrm{end}}}
\newcommand{\Eeasy}{\mathcal E_{\mathrm{easy}}}
\newcommand{\Eres}{\mathcal E_{\mathrm{res}}}
\newcommand{\wt}{\widetilde}
\DeclareMathOperator{\supp}{supp}
\DeclareMathOperator{\TV}{TV}
\DeclareMathOperator{\Res}{Res}

\title{\textbf{Audited SOMK--GFT Reduction for Binary Goldbach}\\
\large Exact Endpoint Reconstruction, Vaughan Raw--Quotient Compilation,\\
and the Structured Outer--Modulus Kloosterman Residual}
\author{Draft manuscript}
\date{20 August 2026}

\begin{document}
\maketitle

\begin{abstract}
We give an audited reduction of the binary Goldbach problem to a structured fixed-shift Kloosterman residual, called global \(\SOMK\).  The note consolidates the previously established generalized-Fourier (GFT) endpoint machinery and adds two exact compiler identities: a Vaughan raw--quotient uncompression and a genuine two-raw refinement of the Type-II term.  The audit also repairs two overstatements that are easy to make in this setting.  First, once the one-large-prime GFT leaf is reconstructed at the prime endpoint, every nonempty friable cofactor vanishes exactly; thus a separate ``friable-core PMKF'' theorem is not a genuine endpoint obstruction.  Second, an integral over the \emph{minor arcs} cannot be converted into a fixed determinant equation by full-circle orthogonality.  Consequently the algebraic MRDET normal form and the lossless minor-arc-to-MRDET interface must be distinguished.

We therefore define global \(\SOMK\) as the conjunction of (i) a pre-Bessel, fixed-\(N\), lossless compiler/interface statement that exhausts all residual endpoint minor-arc leaves by structured outer-modulus Kloosterman kernels, and (ii) a log-power saving for those kernels.  With this precise definition we prove, without any further unproved implication, that
\[
 \boxed{\SOMK\Longrightarrow \text{binary Goldbach for every sufficiently large even integer}.}
\]
The proof uses only the exact GFT endpoint projectors, a standard binary-Goldbach major-arc asymptotic, and the stated \(\SOMK\) hypotheses.  A finite-cell backend audit is included: the Bettin--Chandee chamber is derived with the correct prefactors, the balanced \(17/33\) threshold is recovered for the corresponding trilinear kernel, and the Blomer--Pascadi square-root geometry is identified at the central MRDET cell.  These backend calculations are recorded as local solver information and are not silently promoted to a proof of the global interface.  The final appendices collect the reusable GFT/RDET/MRDET/QRFH toolkit and all type-safety warnings needed for further work.
\end{abstract}

\tableofcontents

\section{Status, theorem-level claim, and audit policy}

The purpose of this manuscript is a reduction theorem, not an unconditional proof of Goldbach.  All statements are labelled according to the following policy.

\begin{center}
\begin{tabular}{>{\raggedright\arraybackslash}p{0.23\textwidth} >{\raggedright\arraybackslash}p{0.66\textwidth}}
\toprule
\textbf{Status} & \textbf{Meaning}\\
\midrule
PROVED & An exact algebraic identity, a standard cited theorem, or a deduction proved in this manuscript.\\
VALID BRIDGE & A type-correct local normal form, valid once its stated hypotheses and variable roles have genuinely been reached.\\
BACKEND PASS & A published estimate closes the displayed local kernel in the stated numerical chamber.  This does \emph{not} by itself prove that every Goldbach minor-arc leaf reaches that kernel.\\
OPEN / HYPOTHESIS & A genuine remaining analytic or interface assertion.\\
REJECTED SHORTCUT & A tempting inference that the audit shows is invalid.\\
\bottomrule
\end{tabular}
\end{center}

The main claim is the following.

\begin{theorem}[Audited SOMK--GFT reduction to binary Goldbach]\label{thm:main}
Assume the global \(\SOMK\) hypothesis of \cref{hyp:global-somk}.  Then, for every sufficiently large even integer \(N\),
\begin{equation}\label{eq:main-asymp}
 R_{\vartheta\vartheta}(N)
 =\SG(N)\Jw N+o(N),
\end{equation}
where \(\SG(N)>0\) is the binary-Goldbach singular series and \(\Jw>0\) is the chosen archimedean weight.  In particular, every sufficiently large even \(N\) is the sum of two primes.
\end{theorem}

The full proof is given in \cref{sec:somk-to-gold}.  An important logical point is that the new route does \emph{not} require proving the stronger frozen-mode hypothesis \(\PMKF\).  Instead, after the GFT scale-count decomposition has isolated the one-hit leaf, we reconstruct that leaf at the prime endpoint and attack the resulting fixed-\(N\) prime--prime minor correlation directly.  Thus the new implication is
\[
 \SOMK
 \Longrightarrow [z]\,\mathcal J_{z,1}(N)=o(N)
 \Longrightarrow \text{Goldbach},
\]
not necessarily \(\SOMK\Rightarrow\PMKF\).

\subsection{Nonclaims}

This manuscript does \emph{not} claim any of the following.
\begin{enumerate}[label=\textup{(N\arabic*)},leftmargin=2.4em]
\item It does not prove global \(\SOMK\).
\item It does not claim that a local Bettin--Chandee, Wright, or Blomer--Pascadi estimate already supplies the lossless global minor-arc interface.
\item It does not convert a minor-arc integral into a Diophantine equality by full-circle orthogonality.
\item It does not identify complement-side GFT factors with prime-side Heath--Brown/Vaughan factors merely because their dyadic lengths agree.
\item It does not count a formal convolution factor \(1\) as a genuine long analytic variable unless a leaf gives it positive length.
\item It does not claim that completion of a short reciprocal variable is automatically lossless.
\item It does not use withdrawn numerical improvements.
\end{enumerate}

\section{Weighted Goldbach target and the standard major arcs}\label{sec:goldbach-target}

Write \(\Lambda\) for the von Mangoldt function and
\[
 \vartheta(n):=(\log n)\1_{\Pp}(n).
\]
Let \(W_1,W_2\in C_c^\infty((0,1))\) be nonnegative and suppose
\begin{equation}\label{eq:Jw}
 \Jw:=\int_0^1 W_1(t)W_2(1-t)\,dt>0.
\end{equation}
For an even integer \(N\), define
\begin{align}
 R_{\Lambda\Lambda}(N)
 &:={\sum}_{m+n=N}\Lambda(m)\Lambda(n)W_1(m/N)W_2(n/N),\label{eq:RLL}\\
 R_{\Lambda\vartheta}(N)
 &:={\sum}_{m+n=N}\Lambda(m)\vartheta(n)W_1(m/N)W_2(n/N),\label{eq:RLtheta}\\
 R_{\vartheta\vartheta}(N)
 &:={\sum}_{m+n=N}\vartheta(m)\vartheta(n)W_1(m/N)W_2(n/N).\label{eq:Rtt}
\end{align}

\begin{lemma}[Prime-power replacement]\label{lem:pp}
Uniformly for \(N\ge3\),
\[
 R_{\Lambda\Lambda}(N)=R_{\vartheta\vartheta}(N)
 +O\!\left(N^{1/2}(\log N)^2\right),
\]
and the same error holds when only one copy of \(\Lambda\) is replaced by \(\vartheta\).
\end{lemma}

\begin{proof}
The difference \(\Lambda-\vartheta\) is supported on proper prime powers \(p^k\), \(k\ge2\).  There are \(O(N^{1/2})\) such integers up to \(N\), and both arithmetic weights are \(O(\log N)\).  Summing trivially gives the stated error.
\end{proof}

Put
\[
 S_\Lambda(\alpha):=\sum_n\Lambda(n)W_1(n/N)e(\alpha n),
 \qquad
 S_\vartheta(\alpha):=\sum_n\vartheta(n)W_2(n/N)e(\alpha n),
\]
where \(e(x)=e^{2\pi i x}\).  Let \(P=(\log N)^B\) with \(B\) sufficiently large and take a standard major-arc set
\[
 \Mcal
 :=\bigcup_{q\le P}\ \bigcup_{\substack{a\bmod q\\(a,q)=1}}
 \left\{\alpha\in\T:\left|\alpha-\frac aq\right|\le\frac{P}{qN}\right\},
 \qquad
 \mcal:=\T\setminus\Mcal.
\]

Define
\begin{equation}\label{eq:sing-series}
 \SG(N)
 :=2C_2\prod_{\substack{p\mid N\\p>2}}\frac{p-1}{p-2},
 \qquad
 C_2:=\prod_{p>2}\left(1-\frac1{(p-1)^2}\right).
\end{equation}
For even \(N\), \(\SG(N)\) is positive and bounded below by an absolute positive constant.

\begin{proposition}[Standard weighted binary-Goldbach major arcs]\label{prop:major}
For the preceding choice of smooth weights and major arcs,
\begin{equation}\label{eq:major}
 \int_{\Mcal}S_\Lambda(\alpha)S_\vartheta(\alpha)e(-N\alpha)\,d\alpha
 =\SG(N)\Jw N+o(N)
\end{equation}
uniformly for sufficiently large even \(N\).
\end{proposition}

\begin{proof}[Status and provenance]
This is the standard smooth binary-Goldbach major-arc calculation.  The local factor at an odd prime \(p\nmid N\) has two forbidden residues, whereas for \(p\mid N\) the two residues coincide; the resulting Euler product is \eqref{eq:sing-series}.  The real-place convolution gives \eqref{eq:Jw}.  No statement about cancellation on the minor arcs is used here.  We treat this classical major-arc asymptotic as an unconditional external input.
\end{proof}

\section{Exact GFT endpoint machinery}\label{sec:gft-endpoint}

\subsection{Squarefree Fourier mode and radial prime extraction}

For \(z\in\mathbb C\), set
\begin{equation}\label{eq:gz}
 g_z(n):=\mu^2(n)z^{\omega(n)}.
\end{equation}
For each fixed \(n\), \(g_z(n)\) is a polynomial in \(z\).

\begin{proposition}[Exact prime coefficient]\label{prop:prime-coeff}
For every \(n\ge2\),
\[
 [z]g_z(n)=\1_{\Pp}(n),
 \qquad
 (\log n)[z]g_z(n)=\vartheta(n).
\]
\end{proposition}

\begin{proof}
If \(n\) is not squarefree then \(g_z(n)=0\).  If \(n\) is squarefree then \(g_z(n)=z^{\omega(n)}\), whose \(z\)-coefficient is nonzero exactly when \(\omega(n)=1\), i.e. exactly when \(n\) is prime.
\end{proof}

\begin{lemma}[Radial reconstruction]\label{lem:radial}
For a polynomial \(F\) and every \(r>0\),
\[
 [z]F(z)=\frac1{2\pi r}\int_0^{2\pi}F(re^{i\theta})e^{-i\theta}\,d\theta,
 \qquad
 |[z]F(z)|\le r^{-1}\sup_{|z|=r}|F(z)|.
\]
\end{lemma}

\begin{proof}
This is Cauchy's coefficient formula.
\end{proof}

The freedom to take \(r<1\) is one of the useful GFT resources: frozen modes then carry the damping factor \(r^{\omega(n)}\).  The endpoint coefficient remains exact.

\subsection{Finite large-factor projectors and one-hit collapse}

Fix \(K\ge2\), put \(Y=X^{1/K}\), and define
\[
 \Omega_{>Y}(n):=\sum_{\substack{p^\nu\parallel n\\p>Y}}\nu.
\]
If \(n\le X\), then \(0\le\Omega_{>Y}(n)\le K-1\), hence
\begin{equation}\label{eq:DFT-hit}
 \1_{\{\Omega_{>Y}(n)=j\}}
 =\frac1K\sum_{\ell=0}^{K-1}
 \exp\!\left(\frac{2\pi i\ell}{K}(\Omega_{>Y}(n)-j)\right).
\end{equation}

For squarefree \(n\le X\), write uniquely
\begin{equation}\label{eq:hit-factor}
 n=a p_1\cdots p_j,
 \qquad
 P^+(a)\le Y<p_1<\cdots<p_j.
\end{equation}
Then
\begin{equation}\label{eq:hit-weight}
 g_z(n)=g_z(a)z^j=z^{j+\omega(a)}.
\end{equation}

\begin{proposition}[Leaf-adaptive endpoint collapse]\label{prop:endpoint-collapse}
The coefficient \([z]g_z(n)\) on a fixed hit leaf \eqref{eq:hit-factor} is nonzero only in the two cases
\begin{enumerate}[label=\textup{(\roman*)}]
\item \(j=0\) and \(a\) is prime, necessarily \(a\le Y\);
\item \(j=1\) and \(a=1\), so \(n=p_1>Y\) is prime.
\end{enumerate}
Every \(j\ge2\) leaf vanishes identically after \([z]\) extraction.
\end{proposition}

\begin{proof}
By \eqref{eq:hit-weight}, the \(z\)-exponent is \(j+\omega(a)\).  Its coefficient of \(z\) is nonzero exactly when \(j+\omega(a)=1\).  The two listed cases are the only nonnegative integer solutions.
\end{proof}

Taking \(K=4\), hence \(Y=N^{1/4}\), the \(j=0\) endpoint contributes at most
\begin{equation}\label{eq:j0small}
 O\!\left(N^{1/4}(\log N)^2\right)=o(N)
\end{equation}
in the weighted Goldbach correlation.  The only target-scale endpoint leaf is therefore the one-hit leaf, and reconstructing that leaf gives the prime-supported sequence on \(p>N^{1/4}\).

\begin{corollary}[No separate friable-core endpoint obstruction]\label{cor:fc-collapse}
Suppose a one-hit squarefree integer is written \(n=dp\), with \(p>Y\) prime and \(P^+(d)\le Y\).  Then
\[
 g_z(dp)=z\,g_z(d)=z^{1+\omega(d)}.
\]
Consequently
\[
 [z]\,g_z(dp)=\1_{\{d=1\}}.
\]
Thus every nonempty cofactor \(d>1\) vanishes \emph{exactly} if the one-hit leaf is reconstructed at the prime endpoint.
\end{corollary}

\begin{proof}
The integer \(d\) is squarefree and coprime to \(p\).  Hence \(g_z(dp)=z^{1+\omega(d)}\).  The coefficient of \(z\) is nonzero if and only if \(\omega(d)=0\), i.e. \(d=1\).
\end{proof}

\begin{remark}[Relation to the older PMKF route]
The older \(\PMKF\) theorem keeps the one-hit leaf frozen in \(z\) and asks for a uniform modewise bound before reconstruction.  That is a stronger sufficient route.  The present SOMK route instead reconstructs the one-hit leaf first and attacks the resulting fixed-\(N\) prime--prime endpoint.  Therefore the new route can imply Goldbach without proving the full frozen-mode \(\PMKF\) hypothesis.
\end{remark}

\section{Exact QRFH refinement}\label{sec:qrfh}

Although QRFH is not needed in the final SOMK implication after \cref{cor:fc-collapse}, it is a valid structural tool and is retained for later use.

Fix \(Y=X^{1/4}\) and choose \(0<\tau<1/4\), \(Z=X^\tau\).  For a one-hit squarefree \(n=dp\le X\) with \(P^+(d)\le Y<p\), factor
\[
 d=bq_1\cdots q_j,
 \qquad
 P^+(b)\le Z<q_1<\cdots<q_j\le Y.
\]

\begin{proposition}[Quotient-retention factor-hit lemma, structural form]\label{prop:qrfh}
The preceding factorization is unique and
\[
 j<\frac{3}{4\tau}.
\]
Moreover
\[
 g_z(n)=z^{j+1}g_z(b).
\]
If a type-correct fixed-shift state has reached
\[
 Ac+dp=N,
\]
then for every \(I\subseteq\{1,\dots,j\}\),
\[
 p\prod_{i\in I}q_i\mid N-Ac.
\]
If additionally \(\bigl(A,p\prod_{i\in I}q_i\bigr)=1\), then
\[
 c\equiv N\overline A\pmod{p\prod_{i\in I}q_i}.
\]
\end{proposition}

\begin{proof}
Uniqueness is unique factorization.  Since \(p>X^{1/4}\) and \(n\le X\), one has \(d<X^{3/4}\).  Each \(q_i>X^\tau\), hence \(X^{j\tau}<q_1\cdots q_j\le d<X^{3/4}\), proving \(j<3/(4\tau)\).  The GFT identity follows from squarefreeness and coprimality.  Finally,
\[
 N-Ac=dp=bpq_1\cdots q_j,
\]
so every displayed modulus divides \(N-Ac\); invert \(A\) modulo it when the stated coprimality holds.
\end{proof}

\section{Exact Vaughan raw--quotient compiler}\label{sec:vaughan}

This section records the main new exact algebraic compiler used by the SOMK route.  The purpose is to avoid the standard loss of structure in which a genuine convolution coefficient is replaced too early by an arbitrary bounded sequence.

For a parameter \(U\ge1\), define
\[
 \mu_{\le U}(n):=\mu(n)\1_{n\le U},
 \qquad
 \mu_{>U}:=\mu-\mu_{\le U},
\]
and similarly \(\Lambda_{\le V}\), \(\Lambda_{>V}\).

\begin{proposition}[Exact Vaughan identity in compiler form]\label{prop:vaughan}
For every \(U,V\ge1\), coefficientwise,
\begin{equation}\label{eq:vaughan}
 \boxed{
 \Lambda
 =\mu_{\le U}*\log
 -\mu_{\le U}*\Lambda_{\le V}*1
 +\Lambda_{\le V}
 +\mu_{>U}*\Lambda_{>V}*1.}
\end{equation}
\end{proposition}

\begin{proof}
Because \(\Lambda=\mu*\log\) and \(\log=\Lambda*1\), we have
\[
 \mu*\Lambda*1=\Lambda,
 \qquad
 \mu*\Lambda_{\le V}*1=\Lambda_{\le V}*(\mu*1)=\Lambda_{\le V}.
\]
Expanding
\[
 \mu_{>U}*\Lambda_{>V}*1
 =(\mu-\mu_{\le U})*(\Lambda-\Lambda_{\le V})*1
\]
gives
\[
 \Lambda-\Lambda_{\le V}-\mu_{\le U}*\log
 +\mu_{\le U}*\Lambda_{\le V}*1.
\]
Rearranging proves \eqref{eq:vaughan}.
\end{proof}

The last term is the classical Type-II term.  Its usual compression obscures a genuine coefficient-one quotient.  The following identity recovers it exactly.

\begin{proposition}[Vaughan raw--quotient uncompression]\label{prop:raw-quotient}
Let
\[
 T_{II}:=\mu_{>U}*\Lambda_{>V}*1.
\]
Then, for every \(n\ge1\),
\begin{equation}\label{eq:raw-quotient}
 \boxed{
 T_{II}(n)
 =-
 \sum_{\substack{m e f=n\\m>V,\ e\le U,\ ef>U}}
 \Lambda(m)\mu(e).}
\end{equation}
In every nonzero summand, \(f>1\).  Thus \(f\) is a genuine coefficient-one quotient, not a formal convolution leg.
\end{proposition}

\begin{proof}
From the definition,
\[
 T_{II}(n)
 =\sum_{\substack{m r=n\\m>V}}
 \Lambda(m)
 \sum_{\substack{a\mid r\\a>U}}\mu(a).
\]
If the inner sum is nonzero then necessarily \(r>U\), hence \(r>1\).  For \(r>1\),
\[
 \sum_{a\mid r}\mu(a)=0,
\]
so
\[
 \sum_{\substack{a\mid r\\a>U}}\mu(a)
 =-\sum_{\substack{e\mid r\\e\le U}}\mu(e).
\]
Writing \(r=ef\) gives \eqref{eq:raw-quotient}, with the necessary condition \(ef=r>U\).  Since \(e\le U\), the inequality \(ef>U\) rules out \(f=1\).
\end{proof}

\begin{proposition}[Genuine two-raw refinement of Type II]\label{prop:two-raw-TII}
The Type-II coefficient also has the exact representation
\begin{equation}\label{eq:two-raw-TII}
 \boxed{
 T_{II}(n)
 =-
 \sum_{\substack{a b e f=n\\ab>V,\ e\le U,\ ef>U}}
 \mu(a)\mu(e)\log b.}
\end{equation}
Every nonzero summand has \(b>1\) and \(f>1\).  Hence \(b\) and \(f\) are two genuine positive-length raw variables after any dyadic leaf on which they have nontrivial scale.
\end{proposition}

\begin{proof}
Insert the exact identity \(\Lambda(m)=(\mu*\log)(m)=\sum_{ab=m}\mu(a)\log b\) into \eqref{eq:raw-quotient}.  The factor \(\log b\) vanishes at \(b=1\), and \cref{prop:raw-quotient} already gives \(f>1\).
\end{proof}

\begin{corollary}[Full-shift Type-II pair determinant state]\label{cor:TT-state}
If two Type-II factors in a \emph{full fixed-shift correlation} are expanded by \cref{prop:two-raw-TII}, every dyadic summand has an exact equation
\begin{equation}\label{eq:TT-MRDET}
 A b f+B d h=N,
 \qquad
 A:=ae,\quad B:=cg,
\end{equation}
where \(b,f,d,h\) are genuine raw variables and \(A,B\) retain the structured M\"obius origins \(\mu(a)\mu(e)\) and \(\mu(c)\mu(g)\).
\end{corollary}

\begin{proof}
Apply \eqref{eq:two-raw-TII} independently to both summands \(n_1,n_2\) in the exact equality \(n_1+n_2=N\), and group \(A=ae\), \(B=cg\).
\end{proof}

\begin{remark}[Smooth localization and cell count]\label{rem:cells}
The conditions \(ab>V\), \(ef>U\), and the size relations inherited from the weights are monomial inequalities.  A standard smooth dyadic partition produces at most \((\log N)^{O(1)}\) leaves.  Such a polylogarithmic cell count can be absorbed by requesting a larger log-power saving from the final uniform solver.  The partition does not license forgetting the factor origins inside \(A\) and \(B\).
\end{remark}

\section{Critical audit: a minor-arc integral is not a full-shift determinant}\label{sec:minor-audit}

This is the most important correction introduced by the present audit.

\begin{lemma}[Restricted orthogonality kernel]\label{lem:restricted-orth}
Let \(E\subseteq\T\) be measurable.  For integers \(m,n,N\),
\[
 \int_E e(\alpha(m+n-N))\,d\alpha
 =\widehat{\1_E}(N-m-n),
\]
where
\[
 \widehat{\1_E}(t):=\int_E e(-t\alpha)\,d\alpha.
\]
In particular, the Kronecker delta \(\1_{m+n=N}\) is obtained only for \(E=\T\).
\end{lemma}

\begin{proof}
This is the definition of the Fourier coefficient of \(\1_E\).  For \(E=\T\), ordinary character orthogonality gives \(\widehat{\1_\T}(t)=\1_{t=0}\).  For a proper minor-arc set \(E=\mcal\), the Fourier coefficients are generally nonzero for \(t\ne0\).
\end{proof}

\begin{warning}[Rejected minor-arc determinant shortcut]\label{warn:minor-det}
One may not replace
\[
 \int_{\mcal}F(\alpha)G(\alpha)e(-N\alpha)\,d\alpha
\]
by a sum restricted to \(m+n=N\) merely by invoking orthogonality.  Any pre-Bessel route from the minor-arc correlation to RDET/MRDET must preserve the Fourier kernel created by the cutoff, or must provide a separate lossless dispersion/localization argument.  The exact identities of \cref{sec:vaughan} remain valid inside the exponential sums, but they do not by themselves prove the global minor-arc-to-MRDET interface.
\end{warning}

This warning is why the final global hypothesis below has two logically distinct components: a compiler/interface component and an analytic Kloosterman component.

\section{Type-correct local RDET/MRDET bridges}\label{sec:mrdet}

The following local bridges are valid once a fixed-shift state has genuinely been isolated.  They are not asserted to arise from \(\mcal\) by \cref{lem:restricted-orth} alone.

\subsection{RDET}

Suppose a routed state has
\begin{equation}\label{eq:RDET}
 uv+dmn=N.
\end{equation}
Let \(g=(u,dm)\).  Then \(g\mid N\).  If \((d,N)=1\), then \(g\mid m\), and with
\[
 u=gq,\qquad m=gm',\qquad R=N/g,
\]
one has
\[
 qv+dm'n=R,
 \qquad
 (q,dm')=1.
\]
After the required gcd extraction, a genuine raw variable can lie in a unique residue class modulo \(q\), and Poisson summation produces reciprocal phases of the shape
\[
 e\!\left(\frac{hR\overline{dn}}q\right).
\]

\subsection{MRDET}

Suppose a type-correct multi-raw leaf has
\begin{equation}\label{eq:MRDET}
 Acv+dBxy=N,
\end{equation}
where \(c,v,x,y\) are genuine raw/smooth coordinates and \(A,B\) carry structured arithmetic coefficients.  Fix \(A,B,c\) and put \(q=Ac\).  Under the required coprimality,
\[
 dBxy\equiv N\pmod q.
\]
Poisson summation in \(x,y\) produces, on the nonzero dual orbit, classical Kloosterman sums
\begin{equation}\label{eq:mrdet-kl}
 S\!\left(Nk\overline{dB},-h;q\right).
\end{equation}
If the raw lengths are \(X_1,X_2\), the corresponding dual lengths are of order \(q/X_1\), \(q/X_2\).

\begin{remark}[Why \cref{cor:TT-state} is useful]
The identity \eqref{eq:TT-MRDET} has exactly the required algebraic shape with \(d=1\):
\[
 A\,b\,f+B\,d\,h=N.
\]
The crucial point is that \(b,f,d,h\) are genuine raw variables supplied by exact identities; they are not manufactured by a formal \(1\)-convolution.  What remains open globally is the lossless passage from the minor-arc correlation to a finite family of such fixed-shift states.
\end{remark}

\section{Definition of the global SOMK--GFT hypothesis}\label{sec:somk-def}

We now define the exact hypothesis used in the main theorem.  The definition is intentionally stronger than a single local Kloosterman estimate and intentionally weaker than an arbitrary-coefficient Goldbach Type-II oracle.

\subsection{Endpoint minor correlation}

Let
\begin{equation}\label{eq:Eend}
 \Eend(N)
 :=\int_{\mcal}S_\Lambda(\alpha)S_\vartheta(\alpha)e(-N\alpha)\,d\alpha.
\end{equation}
This is the only target-scale endpoint minor contribution after the exact GFT reconstruction of \cref{prop:endpoint-collapse}, up to the negligible small-prime leaf \eqref{eq:j0small}.

\begin{lemma}[Minor-arc replacement of \(\vartheta\) by \(\Lambda\)]\label{lem:minor-pp}
Let
\[
 D(\alpha):=\sum_n(\Lambda(n)-\vartheta(n))W_2(n/N)e(\alpha n).
\]
Then
\begin{equation}\label{eq:minor-pp}
 \int_{\mcal}S_\Lambda(\alpha)D(\alpha)e(-N\alpha)\,d\alpha
 \ll N^{3/4}(\log N)^{3/2}.
\end{equation}
In particular, the difference is \(o(N)\).
\end{lemma}

\begin{proof}
By Cauchy--Schwarz and Parseval,
\[
 \left|\int_{\mcal}S_\Lambda D e(-N\alpha)\,d\alpha\right|
 \le \|S_\Lambda\|_2\|D\|_2.
\]
The standard second moment gives
\[
 \|S_\Lambda\|_2^2
 =\sum_n\Lambda(n)^2W_1(n/N)^2
 \ll N\log N.
\]
The sequence \(\Lambda-\vartheta\) is supported on proper prime powers.  There are \(O(N^{1/2})\) of them and each coefficient is \(O(\log N)\), hence
\[
 \|D\|_2^2\ll N^{1/2}(\log N)^2.
\]
Multiplying the square roots proves \eqref{eq:minor-pp}.
\end{proof}

Thus the endpoint minor problem can be treated, at a cost \(o(N)\), as a \(\Lambda\times\Lambda\) problem, on which \cref{prop:vaughan,prop:raw-quotient,prop:two-raw-TII} may be applied on both sides.

\subsection{Local structured outer-modulus Kloosterman kernel}

A \emph{local SOMK kernel} is an actual contribution, including all Poisson prefactors and smooth weights, of the schematic form
\begin{equation}\label{eq:local-somk}
 \Kl_\rho(N)
 =\Xi_\rho
 \sum_{q\sim Q}\gamma_{\rho,u}(q)
 \sum_{B\sim B_0}\beta_{\rho,v}(B)
 \sum_{k\sim K}\alpha_{\rho}(k)
 \sum_{\ell\sim L}\nu_{\rho}(\ell)
 S\!\left(N_\rho k\overline B,-\ell;q\right),
\end{equation}
where:
\begin{enumerate}[label=\textup{(K\arabic*)},leftmargin=2.7em]
\item \(N_\rho=N/g\) with \(g\mid N\) arising from a legitimate gcd extraction;
\item the inverse \(\overline B\pmod q\) is used only on branches where \((B,q)=1\);
\item \(\Xi_\rho\) is the exact normalization/prefactor produced by the compiler and is part of the kernel, so no scale factor is silently dropped;
\item \(\gamma_{\rho,u}\) and \(\beta_{\rho,v}\) retain certified factor origins (for example M\"obius factors, prime factors, GFT modes, and factor-scale tags) until a downstream theorem explicitly permits their replacement by arbitrary coefficients;
\item the total \(L^2\) norms and smooth-derivative costs are bounded by fixed powers of \(\log N\) times the natural square-root lengths;
\item the auxiliary GFT/Wiener frequencies \(u,v\) are retained, with at most polynomial frequency loss.
\end{enumerate}
The schematic display \eqref{eq:local-somk} is not itself a theorem; the point of the interface hypothesis below is to assert that the actual residual endpoint minor leaves reach such kernels losslessly.

\subsection{The two components}

\begin{hypothesis}[Pre-Bessel MRDET interface, \(\PBM\)]\label{hyp:pbm}
For every \(A>0\), after:
\begin{enumerate}[label=\textup{(I\arabic*)},leftmargin=2.7em]
\item exact GFT endpoint reconstruction;
\item the minor-arc prime-power replacement of \cref{lem:minor-pp};
\item the exact Vaughan decomposition \eqref{eq:vaughan};
\item raw--quotient uncompression \eqref{eq:raw-quotient} and, on residual Type-II leaves, the genuine two-raw refinement \eqref{eq:two-raw-TII};
\item smooth dyadic/factor-scale localization with type labels retained;
\end{enumerate}
the endpoint minor correlation has a decomposition
\begin{equation}\label{eq:pbm-decomp}
 \Eend(N)
 =\Eeasy(N)
 +\sum_{\rho\in\mathscr R_N}c_\rho\Kl_\rho(N)
 +O_A\!\left(\frac{N}{(\log N)^A}\right),
\end{equation}
where
\[
 |\Eeasy(N)|\ll_A\frac{N}{(\log N)^A},
 \qquad
 |\mathscr R_N|+\sum_{\rho\in\mathscr R_N}|c_\rho|
 \ll (\log N)^{C_0}
\]
for some fixed \(C_0\), and every \(\Kl_\rho\) is a type-correct local SOMK kernel in the sense of \eqref{eq:local-somk}.

The decomposition must be established without averaging over \(N\) and without using the invalid shortcut of \cref{warn:minor-det}.
\end{hypothesis}

\begin{hypothesis}[Local SOMK saving]\label{hyp:local-somk}
For every \(B>0\), uniformly for every sufficiently large fixed even \(N\), every admissible residual leaf \(\rho\in\mathscr R_N\), and every admissible auxiliary GFT/Wiener mode,
\begin{equation}\label{eq:local-somk-bound}
 \boxed{
 |\Kl_\rho(N)|
 \ll_{B}\frac{N}{(\log N)^B}}
\end{equation}
with at most polynomial loss in the auxiliary Fourier frequencies.  The polynomial frequency loss is required to be compatible with the weighted Wiener norm of the factor-scale reconstruction.
\end{hypothesis}

\begin{hypothesis}[Global SOMK--GFT]\label{hyp:global-somk}
The conjunction of \cref{hyp:pbm} and \cref{hyp:local-somk} holds.
\end{hypothesis}

\begin{remark}[Why the interface is explicitly part of global SOMK]
The algebraic identities of \cref{sec:vaughan,sec:mrdet} show that the desired raw variables and determinant/Kloosterman geometry exist once an appropriate fixed-shift leaf has been isolated.  They do \emph{not} by themselves prove that the minor-arc cutoff reaches this geometry without loss.  Packaging the lossless interface as an explicit component of global SOMK prevents the main theorem from hiding an unproved arrow.
\end{remark}

\section{Strict proof that global SOMK--GFT implies Goldbach}\label{sec:somk-to-gold}

We now prove \cref{thm:main}.

\begin{proposition}[Global SOMK gives the endpoint minor-arc saving]\label{prop:somk-minor}
Assume global \(\SOMK\).  Then for every \(A>0\),
\begin{equation}\label{eq:Eend-small}
 \Eend(N)\ll_A\frac{N}{(\log N)^A}
\end{equation}
for every sufficiently large even \(N\).
\end{proposition}

\begin{proof}
Fix \(A>0\).  Apply \cref{hyp:pbm} with exponent \(A+2C_0+10\).  This gives \eqref{eq:pbm-decomp}, with the easy term and compiler error bounded by \(N(\log N)^{-A-2C_0-10}\).

Apply \cref{hyp:local-somk} to each residual kernel with exponent \(B=A+2C_0+10\).  Then
\[
 \left|\sum_{\rho\in\mathscr R_N}c_\rho\Kl_\rho(N)\right|
 \le
 \left(\sum_{\rho\in\mathscr R_N}|c_\rho|\right)
 \sup_{\rho\in\mathscr R_N}|\Kl_\rho(N)|
 \ll
 (\log N)^{C_0}\frac{N}{(\log N)^{A+2C_0+10}}.
\]
This is \(O_A(N(\log N)^{-A})\).  The weighted Wiener reconstruction contributes at most a further fixed/polylogarithmic factor by hypothesis and is absorbed by the same excess exponent.  Combining all terms proves \eqref{eq:Eend-small}.
\end{proof}

\begin{proposition}[Complete minor arcs are negligible]\label{prop:all-minor}
Assume global \(\SOMK\).  Then
\begin{equation}\label{eq:minor-all}
 [z]\int_{\mcal}S_\Lambda(\alpha)S_z(\alpha)e(-N\alpha)\,d\alpha=o(N),
\end{equation}
where
\[
 S_z(\alpha):=\sum_n(\log n)g_z(n)W_2(n/N)e(\alpha n).
\]
Equivalently,
\[
 \int_{\mcal}S_\Lambda(\alpha)S_\vartheta(\alpha)e(-N\alpha)\,d\alpha=o(N).
\]
\end{proposition}

\begin{proof}
Use \eqref{eq:DFT-hit} with \(K=4\), \(Y=N^{1/4}\), and reconstruct each hit leaf separately.  By \cref{prop:endpoint-collapse}, the \(j=2,3\) leaves vanish exactly after \([z]\); the \(j=0\) leaf contributes \(O(N^{1/4}(\log N)^2)=o(N)\) by \eqref{eq:j0small}; and the \(j=1\) leaf reconstructs exactly to the prime-supported endpoint with prime \(>N^{1/4}\).  Its minor-arc contribution is precisely \(\Eend(N)\), up to the already counted small-prime leaf.  By \cref{prop:somk-minor}, \(\Eend(N)=o(N)\).  This proves \eqref{eq:minor-all}.  The equivalent \(S_\vartheta\) formulation follows from \cref{prop:prime-coeff} and linearity.
\end{proof}

\begin{proof}[Proof of \cref{thm:main}]
Split the exact circle identity
\[
 R_{\Lambda\vartheta}(N)
 =\int_\T S_\Lambda(\alpha)S_\vartheta(\alpha)e(-N\alpha)\,d\alpha
\]
into major and minor arcs.  The major arcs equal \(\SG(N)\Jw N+o(N)\) by \cref{prop:major}.  The minor arcs are \(o(N)\) by \cref{prop:all-minor}.  Hence
\begin{equation}\label{eq:RLtheta-asymp}
 R_{\Lambda\vartheta}(N)=\SG(N)\Jw N+o(N).
\end{equation}
Since \(\SG(N)\) and \(\Jw\) are positive, the right side is positive for all sufficiently large even \(N\).

Finally, replace the remaining \(\Lambda\) by \(\vartheta\).  By \cref{lem:pp},
\[
 R_{\vartheta\vartheta}(N)
 =R_{\Lambda\vartheta}(N)
 +O\!\left(N^{1/2}(\log N)^2\right)
 =\SG(N)\Jw N+o(N).
\]
Thus \(R_{\vartheta\vartheta}(N)>0\) for every sufficiently large even \(N\).  Every nonzero term in \(R_{\vartheta\vartheta}(N)\) is a representation \(N=p_1+p_2\) by two primes.  This proves binary Goldbach for all sufficiently large even \(N\).
\end{proof}

\begin{corollary}[Full Goldbach after finite verification]\label{cor:finite}
If global \(\SOMK\) is proved with an explicit threshold \(N\ge N_0\), and binary Goldbach is verified directly for the even integers \(4\le N<N_0\), then every even integer greater than \(2\) is the sum of two primes.
\end{corollary}

\section{Finite-cell backend audit}\label{sec:cell-audit}

This section records local solver information that becomes available \emph{after} a type-correct pre-Bessel interface has reached the displayed kernel.  None of these calculations is used to smuggle in \cref{hyp:pbm}.

\subsection{A three-variable reciprocal kernel and its natural dual length}

Consider a legitimate smooth fixed-shift state
\begin{equation}\label{eq:qfrh}
 qf+rh=N,
\end{equation}
with
\[
 q\asymp Q=X^\alpha,
 \quad
 r\asymp R=X^\beta,
 \quad
 f\asymp F=X/Q,
 \quad
 h\asymp H=X/R.
\]
Assume \((q,r)=1\).  Summing over \(h\equiv N\overline r\pmod q\) and applying Poisson gives a nonzero-frequency kernel with natural dual length
\begin{equation}\label{eq:Kdual}
 K\asymp\frac qH\asymp\frac{QR}{X}=X^\kappa,
 \qquad
 \kappa:=\alpha+\beta-1.
\end{equation}
The Poisson prefactor is
\[
 \frac Hq\asymp X^{-\kappa}.
\]
The phase is \(e_q(kN\overline r)\).  The conductor ratio appearing in Bettin--Chandee is
\begin{equation}\label{eq:BCconductor}
 \frac{N K}{RQ}\asymp1,
\end{equation}
so the true Goldbach homogeneity eliminates the spurious fixed positive ``additive defect'' loss that appears in an over-large outer envelope.

If \(\kappa<0\) by a fixed amount, then \(K<1\) and every nonzero Fourier mode is power-negligible by smooth Fourier decay.  This observation concerns the local fixed-shift kernel; it does not, by itself, replace \cref{hyp:pbm} for a minor-arc integral.

\subsection{Bettin--Chandee chamber}

We use the Bettin--Chandee trilinear estimate in the form
\[
 \sum_{a,m,n}\nu_a\alpha_m\beta_n
 e\!\left(\vartheta\frac{a\overline m}{n}\right)
\]
with arbitrary coefficients, together with their published bound \cite{BC}.  Map
\[
 a\leftrightarrow k\ (K),
 \qquad
 m\leftrightarrow r\ (R),
 \qquad
 n\leftrightarrow q\ (Q),
 \qquad
 \vartheta\leftrightarrow N.
\]
Assume the coefficient \(L^2\)-norms have their natural square-root sizes up to \((\log X)^{O(1)}\).

\begin{proposition}[Audited BC local cell guard]\label{prop:bc-guard}
Let \(m_*:=\max(\alpha,\beta)\) and \(\kappa=\alpha+\beta-1\ge0\).  After restoring the Poisson prefactor \(H/q\), the two Bettin--Chandee terms have exponents
\begin{align}
 E_1&=\frac{17}{20}+\frac7{10}\kappa+\frac14m_*,\label{eq:E1}\\
 E_2&=\frac78+\frac78\kappa+\frac18m_*\label{eq:E2}
\end{align}
relative to \(X\).  Hence the local kernel has a power saving over the target scale \(X\) whenever
\begin{align}
 14\kappa+5m_*&<3,\label{eq:BC1}\\
 7\kappa+m_*&<1.\label{eq:BC2}
\end{align}
Equivalently, the first guard is
\begin{equation}\label{eq:BC1-alt}
 14(\alpha+\beta)+5\max(\alpha,\beta)<17.
\end{equation}
Inside the feasible region of \eqref{eq:BC1}, \eqref{eq:BC2} is automatic.  On the diagonal \(\alpha=\beta=t\), the frontier is
\begin{equation}\label{eq:17-33}
 t<\frac{17}{33}.
\end{equation}
\end{proposition}

\begin{proof}
The product of the three coefficient norms is
\[
 (KRQ)^{1/2}=X^{1/2+\kappa}.
\]
Multiplying by the Poisson prefactor \(X^{-\kappa}\) leaves \(X^{1/2}\).  By \eqref{eq:BCconductor}, the Bettin--Chandee conductor factor is \(O(1)\).

For their first term,
\[
 (KRQ)^{7/20}(R+Q)^{1/4}
 \asymp
 X^{\frac7{20}(1+2\kappa)+m_*/4}.
\]
Including the surviving \(X^{1/2}\) gives \eqref{eq:E1}.  The condition \(E_1<1\) is \eqref{eq:BC1}.

For the second term,
\[
 (KRQ)^{3/8}(KQ+KR)^{1/8}
 \asymp
 X^{\frac38(1+2\kappa)+(\kappa+m_*)/8}.
\]
Including \(X^{1/2}\) gives \eqref{eq:E2}, and \(E_2<1\) is \eqref{eq:BC2}.

To see that \eqref{eq:BC2} is automatic in the feasible \eqref{eq:BC1} region, note that \(m_*\ge(1+\kappa)/2\).  Feasibility of \eqref{eq:BC1} therefore forces
\[
 14\kappa+\frac52(1+\kappa)<3,
\]
so \(\kappa<1/33\).  Combining this with \(m_*<(3-14\kappa)/5\) yields \(7\kappa+m_*<1\).  Finally, on \(\alpha=\beta=t\), \eqref{eq:BC1-alt} becomes \(33t<17\).
\end{proof}

\begin{remark}[A genuine BC residual edge]
The cell \((\alpha,\beta)=(2/3,1/3)\) has \(\kappa=0\) but violates \eqref{eq:BC1}.  Thus the BC backend does not close the natural \(2/3\)-by-\(1/3\) boundary cell.
\end{remark}

\subsection{Blomer--Pascadi central MRDET geometry}

At a balanced genuine MRDET leaf one may have
\[
 A,B,b,f,d,h\asymp X^{1/3}
\]
in the sense that each side has three multiplicative scale contributions summing to exponent \(1\).  If the modulus is \(q=Ab\asymp X^{2/3}\) and Poisson is applied in raw variables of length \(X^{1/3}\), then the dual lengths satisfy
\[
 K,L\asymp X^{1/3}=q^{1/2}.
\]
Blomer--Pascadi's 2026 theorem gives a \(q^{-1/32+o(1)}\) saving for a fixed modulus in this square-root range \cite{BP}.  This is an exact geometric match at the level of the inner bilinear Kloosterman form.

\begin{warning}[Outer-modulus obstruction]
The local Blomer--Pascadi saving is not automatically preserved after summing absolutely over the structured outer variables \(q\) and \(B\).  The remaining analytic challenge is precisely to exploit their retained M\"obius/GFT/factor-origin structure rather than erase it before applying the fixed-modulus estimate.
\end{warning}

\subsection{Wright and partially fixed moduli}

Wright's current v2 (revised 7 August 2026) improves trilinear Kloosterman-fraction estimates when the denominator has a fixed factor and derives new unbalanced convolution ranges \cite{Wright}.  This is compatible with leaves in which a factor origin inside the modulus is retained.  We record it as a candidate backend, not as a global solver: a complete application must retain the theorem's precise coefficient-distribution hypotheses and all outer prefactors.

\section{Where almost-all Goldbach loses fixed-\(N\) information}\label{sec:akeno}

Akeno proves that for almost all even \(N\), the Goldbach-prime set has level of distribution \(1/6\) \cite{Akeno}.  The point relevant to the present project is not merely the final exponent but the exact place at which fixed-\(N\) information is averaged away.

In Akeno's Lemma 3.12, the minor-arc quantity is squared and summed over \(X/2<N\le X\); the proof then explicitly invokes Cauchy--Schwarz and Bessel's inequality to bound that \(N\)-average.  Consequently the result is an \(\ell^2_N\) statement from which an exceptional set is removed.  It does not supply a pointwise-in-\(N\) estimate for the individual Fourier coefficient.

\begin{proposition}[Wiener reconstruction does not upgrade \(\ell^2_N\) to \(\ell^\infty_N\)]\label{prop:no-l2-linf}
Suppose a family of modes \(E_u(N)\) satisfies
\[
 \sum_N|E_u(N)|^2\le B(u)^2.
\]
If \(E(N)=\int E_u(N)\,d\nu(u)\), then Minkowski gives only
\[
 \|E\|_{\ell^2_N}
 \le\int B(u)\,d|\nu|(u).
\]
No pointwise-in-\(N\) conclusion follows without additional structure.
\end{proposition}

\begin{proof}
Minkowski's integral inequality gives the displayed \(\ell^2\) bound.  For the failure of an automatic \(\ell^\infty\) upgrade, take all modes to be concentrated on the same single integer \(N_0\).  The \(\ell^2\) information then permits the entire norm to sit at \(N_0\).
\end{proof}

Thus global \(\SOMK\) is designed as a genuinely pre-Bessel, fixed-\(N\) replacement for the step that is averaged in Akeno's almost-all theorem.

\section{Current obstacle map}\label{sec:status-map}

\begin{center}
\begin{longtable}{>{\raggedright\arraybackslash}p{0.36\textwidth} >{\raggedright\arraybackslash}p{0.16\textwidth} >{\raggedright\arraybackslash}p{0.38\textwidth}}
\toprule
\textbf{Object} & \textbf{Status} & \textbf{Audit conclusion}\\
\midrule
\endfirsthead
\toprule
\textbf{Object} & \textbf{Status} & \textbf{Audit conclusion}\\
\midrule
\endhead
Radial prime extraction & PROVED & Exact coefficient reconstruction on any fixed circle \(|z|=r\).\\
Finite scale-count DFT & PROVED & Exact projector for the number of factors above \(X^{1/K}\).\\
Multi-hit endpoint leaves & CLOSED & \(j\ge2\) vanish exactly after leaf-local \([z]\) extraction.\\
Zero-hit endpoint & CLOSED & Supported on primes \(\le X^{1/4}\), hence \(o(N)\).\\
One-hit nonempty cofactor & CLOSED at endpoint & \([z]\,z g_z(d)=0\) for every squarefree \(d>1\).\\
QRFH structural second hit & PROVED & Genuine factors of the one-hit cofactor give genuine moduli in a type-correct fixed-shift state.\\
Vaughan raw quotient & PROVED & Type II contains a genuine coefficient-one quotient \(f>1\).\\
Two-raw Type-II refinement & PROVED & After \(\Lambda=\mu*\log\), the raw legs \(b,f\) are both genuine.\\
Full-shift TT\(\times\)TT MRDET algebra & PROVED & Exact equation \(Abf+Bdh=N\) in a full fixed-shift correlation.\\
Minor-arc \(\Rightarrow\) determinant by orthogonality & REJECTED & Minor-arc cutoff leaves a Fourier kernel; full-circle orthogonality cannot be used.\\
Global pre-Bessel minor-arc \(\Rightarrow\) SOMK interface & OPEN & This is \cref{hyp:pbm}.\\
BC local chamber & BACKEND PASS & Guard \(14(\alpha+\beta)+5\max(\alpha,\beta)<17\); diagonal frontier \(17/33\).\\
BP central fixed-modulus kernel & BACKEND PASS locally & Square-root geometry matches; outer structured modulus averaging remains open.\\
Local/global SOMK saving & OPEN & This is \cref{hyp:local-somk}.\\
Global \(\SOMK\Rightarrow\) Goldbach & PROVED & \cref{thm:main}.\\
\bottomrule
\end{longtable}
\end{center}

\section{Research target after the audit}\label{sec:research-target}

The remaining problem has two sharply separated parts.

\paragraph{Interface problem.}
Prove \(\PBM\): retain the minor-arc cutoff and the fixed value of \(N\), perform the Vaughan/GFT raw compiler without a global Bessel average, and derive a finite/polylogarithmic family of actual SOMK kernels with controlled errors.

\paragraph{Analytic problem.}
Prove the local SOMK saving \eqref{eq:local-somk-bound}, exploiting the fact that the outer modulus coefficients are not arbitrary.  At the central MRDET cell the inner dual variables are in the Blomer--Pascadi square-root regime, so the most promising route is to keep the outer factor origins visible while applying the fixed-modulus bilinear theorem inside.  Other cells may be routed to Bettin--Chandee or Wright according to their certified geometry.

If both problems are solved, no additional parity theorem is required by the present reduction: \cref{thm:main} then gives binary Goldbach for all sufficiently large even \(N\).

\appendix

\section{Reusable GFT toolkit}\label{app:toolkit}

This appendix collects exact tools that are safe to reuse.  Some are not needed in the short proof of \cref{thm:main}, but they preserve structure for future SOMK attacks.

\subsection{Tensor factor-scale synthesis and weighted Wiener closure}

Let \(\phi_1,\dots,\phi_d\) be real-valued observables on primes and define
\[
 A_j(n):=\sum_{p\mid n}\phi_j(p),
 \qquad
 \bm A(n):=(A_1(n),\dots,A_d(n)).
\]
For \(u=(\theta,\bm t)\in\T\times\mathbb R^d\), put
\[
 G_u(n)
 :=\mu^2(n)\exp\!\left(i\theta\omega(n)+i\bm t\cdot\bm A(n)\right).
\]
For \((m,n)=1\), \(G_u(mn)=G_u(m)G_u(n)\); for arbitrary \(m,n\), the exact identity includes \(\1_{(m,n)=1}\).

Let \(\Psi\in\mathcal S(\mathbb R^d)\) and define
\[
 d\nu_{k,\Psi}
 :=\frac{e^{-ik\theta}}{2\pi}
 \frac{\widehat\Psi(\bm t)}{(2\pi)^d}\,d\theta\,d\bm t.
\]
Then
\[
 \int G_u(n)\,d\nu_{k,\Psi}(u)
 =\mu^2(n)\1_{\{\omega(n)=k\}}\Psi(\bm A(n)).
\]
For
\[
 W_s(\Psi):=\frac1{(2\pi)^d}\int_{\mathbb R^d}
 |\widehat\Psi(\bm t)|(1+\|\bm t\|)^s\,d\bm t,
\]
any modewise bound
\[
 |\mathcal J_{\theta,\bm t}|
 \le M(1+\|\bm t\|)^s
\]
transfers as
\[
 \left|\int\mathcal J_u\,d\nu_{k,\Psi}(u)\right|
 \le M W_s(\Psi).
\]
This is an \(L^1\) Fourier/Wiener transference and introduces no additive-frequency Bessel diagonal by itself.

\subsection{Complex charge--dimension transport}

Let \(d_z(n)\) be the coefficients of \(\zeta(s)^z\).  Then
\[
 d_{z_1}*d_{z_2}=d_{z_1+z_2},
 \qquad
 d_z=1^{*J}*d_{z-J}\quad(J\ge1).
\]
This is an exact complex convolution identity.  A formal \(1\)-leg is not a genuine long variable unless the leaf gives it positive length.

\subsection{Squarefree one-raw gauge}

There is a multiplicative \(r_z\) such that
\[
 g_z=1*r_z,
\]
with
\[
 r_z(p)=z-1,
 \qquad
 r_z(p^2)=-z,
 \qquad
 r_z(p^k)=0\quad(k\ge3).
\]
Indeed, the local Euler polynomial is
\[
 (1+zx)(1-x)=1+(z-1)x-zx^2.
\]

\subsection{Friable zeta gauge}

For \(Y\ge2\), define
\[
 D_{z,Y}(s):=\prod_{p\le Y}(1-p^{-s})^{-z}
 =\sum_n\frac{d_z^{(\le Y)}(n)}{n^s}.
\]
Then
\[
 D_{z,Y}(s)=\zeta(s)H_{z,Y}(s),
\]
where
\[
 H_{z,Y}(s)
 =\prod_{p\le Y}(1-p^{-s})^{1-z}
 \prod_{p>Y}(1-p^{-s}).
\]
Consequently
\[
 d_z^{(\le Y)}=1*h_{z,Y}.
\]
More generally,
\[
 d_z^{(\le Y)}=1^{*J}*h^{(J)}_{z,Y}.
\]
At first coefficient,
\[
 [z]d_z^{(\le Y)}(n)
 =\begin{cases}
 1/\nu,&n=p^\nu,\ p\le Y,\\
 0,&\text{otherwise}.
 \end{cases}
\]
This illustrates why leaf-local reconstruction can collapse a seemingly highly composite friable leaf to prime-power support.

\subsection{Logarithmic rawization}

For generalized divisor coefficients,
\[
 (\log n)d_z(n)=z(\Lambda*d_z)(n).
\]
For the squarefree GFT and squarefree \(n>1\),
\[
 g_z(n)\log n
 =z\sum_{p\mid n}(\log p)g_z(n/p).
\]
A smooth partition in \(\log p/\log X\) extracts a genuine prime-scale coordinate without erasing the complex GFT mode on the cofactor.

\subsection{Balanced two-raw finite gauge}

Let
\[
 M_U(s):=\sum_{m\le U}\frac{\mu(m)}{m^s},
 \qquad
 \mathcal A_U(s):=\zeta(s)M_U(s),
 \qquad
 \mathcal R_U(s):=1-\mathcal A_U(s).
\]
The coefficients of \(\mathcal R_U\) vanish for \(n\le U\), hence \(\rho_U^{*K}(n)=0\) for \(n\le U^K\).  Since
\[
 (1-R)^2\sum_{j=0}^{K-1}(j+1)R^j
 =1-R^K((K+1)-KR),
\]
one obtains coefficientwise for \(n\le U^K\),
\[
 1=\mathcal A_U^2\sum_{j=0}^{K-1}(j+1)\mathcal R_U^j.
\]
Thus both generalized-divisor and \(\Lambda\)-type coefficients admit exact finite decompositions with two formal raw legs.  This gauge is useful, but unlike \cref{prop:two-raw-TII}, the formal legs need not both have positive length on every leaf.

\subsection{Radial damping}

On \(|z|=r<1\), a squarefree cofactor with \(q\) unresolved prime factors has weight \(r^q\).  Hence
\[
 q\ge C\log\log X
 \quad\Longrightarrow\quad
 r^q\le(\log X)^{-C|\log r|}.
\]
This is a compiler resource for high-factor-count tails, not a substitute for a low-factor analytic estimate.

\section{Type-safety and audit warnings}\label{app:warnings}

\begin{warning}[Variable origin is part of the type]
A factor living on the shifted complement is not the same object as a factor generated by a Heath--Brown or Vaughan decomposition of the prime variable, even when their lengths coincide.  Any cast between them requires an algebraic bridge.
\end{warning}

\begin{warning}[Difference lattices do not automatically save a spectral diagonal]
A condition \(D\mid(n_1-n_2)\) contains the full diagonal \(n_1=n_2\).  Sparsity of off-diagonal pairs therefore does not by itself produce a \(D^{-1}\) spectral saving.
\end{warning}

\begin{warning}[Modulus charge is not automatically an automorphic \(L\)-charge]
An arithmetic weight attached to a modulus/cusp parameter does not automatically become a power of an automorphic \(L\)-function after Kuznetsov.  The corresponding Hecke/local operator must actually be present in the trace formula.
\end{warning}

\begin{warning}[Completion is not automatically lossless]
Completing a short inverse-phase variable can enlarge a length \(M\ll q\) to a full Fourier variable modulo \(q\).  A strong fixed-modulus theorem applied after such completion may fail to recover the geometry that was destroyed.
\end{warning}

\begin{warning}[Formal raw mass is not analytic raw mass]
An identity such as \(d_z=1^{*J}*d_{z-J}\) inserts formal convolution coordinates.  One may count a coordinate as a genuine raw analytic variable only after the localized leaf gives it positive length.  The genuine raw variables in \cref{prop:raw-quotient,prop:two-raw-TII} are stronger because their nontriviality follows from the exact support constraints \(ef>U\) and \(\log b\ne0\).
\end{warning}

\begin{warning}[Minor-arc orthogonality]
The full-circle identity \(\int_\T e(\alpha t)d\alpha=\1_{t=0}\) must never be substituted into an integral over \(\mcal\).  This is the principal correction separating the proved algebraic MRDET compiler from the open \(\PBM\) interface.
\end{warning}

\section{Comparison with the older PMKF--GFT reduction}\label{app:pmkf}

For reference, the older sufficient hypothesis is
\[
 \sup_{|z|=r}|\mathcal J_{z,1}(N)|
 \ll_A \frac{N}{(\log N)^A},
\]
where \(\mathcal J_{z,1}\) is the frozen-mode one-hit minor correlation.  Together with the exact GFT scale split and the standard major arcs, this implies Goldbach for all sufficiently large even \(N\).  That implication remains valid.

The SOMK route is logically different.  It uses the same exact GFT projectors to reach the one-hit leaf, but then reconstructs the leaf at the prime endpoint.  The nonempty one-hit cofactor disappears by \cref{cor:fc-collapse}, and the remaining analytic target is a fixed-\(N\) prime--prime minor correlation.  Thus:
\[
 \begin{array}{c}
 \text{older route: frozen one-hit mode}\ \xrightarrow{\PMKF}\ \text{Goldbach},\\[1mm]
 \text{new route: reconstruct one-hit endpoint}\ \xrightarrow{\SOMK}\ \text{Goldbach}.
 \end{array}
\]
Neither route is asserted to imply the other.

\section{Dependency graph}\label{app:graph}

The audited dependency graph for the new route is
\[
\boxed{
\begin{array}{c}
\text{exact GFT prime projector + scale-count DFT}\\
\Downarrow\\
\text{leaf-adaptive endpoint collapse}\\
\Downarrow\\
\text{prime--prime endpoint minor correlation}\\
\Downarrow\quad\text{(minor prime-power replacement)}\\
\Lambda\times\Lambda\text{ endpoint minor correlation}\\
\Downarrow\quad\text{(exact Vaughan/raw-quotient identities)}\\
\text{structured raw leaf family}\\
\Downarrow\quad\boxed{\PBM\ \text{(OPEN interface)}}\\
\text{finite/polylog family of SOMK kernels}\\
\Downarrow\quad\boxed{\text{local SOMK saving (OPEN)}}\\
\text{endpoint minor arcs}=o(N)\\
\Downarrow\quad\text{(standard major arcs)}\\
R_{\Lambda\vartheta}(N)=\SG(N)\Jw N+o(N)\\
\Downarrow\quad\text{(prime-power replacement)}\\
R_{\vartheta\vartheta}(N)>0\\
\Downarrow\\
\textbf{binary Goldbach for all sufficiently large even }N.
\end{array}}
\]
Every arrow outside the two boxed OPEN inputs is proved in this manuscript or is a standard cited major-arc theorem.

\begin{thebibliography}{99}

\bibitem{Akeno}
M.~Akeno,
\emph{On the level of distribution of Goldbach primes and its applications},
arXiv:2606.29559, v1, 28 June 2026.

\bibitem{ABL}
E.~Assing, V.~Blomer and J.~Li,
\emph{Uniform Titchmarsh divisor problems},
Adv. Math. \textbf{393} (2021), Paper No. 108076;
arXiv:2005.13915.

\bibitem{BC}
S.~Bettin and V.~Chandee,
\emph{Trilinear forms with Kloosterman fractions},
Adv. Math. \textbf{328} (2018), 1234--1262;
arXiv:1502.00769.

\bibitem{BP}
V.~Blomer and A.~Pascadi,
\emph{Bilinear forms with Kloosterman sums via quadratic characters},
arXiv:2607.24311, 27 July 2026.

\bibitem{Chen}
J.-R.~Chen,
\emph{On the representation of a large even integer as the sum of a prime and the product of at most two primes},
Sci. Sinica \textbf{16} (1973), 157--176.

\bibitem{DT}
S.~Drappeau and B.~Topacogullari,
\emph{Combinatorial identities and Titchmarsh's divisor problem for multiplicative functions},
Algebra Number Theory \textbf{13} (2019), 2383--2425;
arXiv:1807.09569.

\bibitem{Pascadi}
A.~Pascadi,
\emph{On the exponents of distribution of primes and smooth numbers},
arXiv:2505.00653, 2025.

\bibitem{Vaughan}
R.~C.~Vaughan,
\emph{Sommes trigonom\'etriques sur les nombres premiers},
C. R. Acad. Sci. Paris S\'er. A-B \textbf{285} (1977), A981--A983.

\bibitem{Wright}
T.~Wright,
\emph{Trilinear Kloosterman fractions I: partially fixed moduli and unbalanced convolutions},
arXiv:2604.25177, v2, 7 August 2026.

\end{thebibliography}

\end{document}
